\chapter{La donnée manquante comme livrable, hiérarchisé}
\label{ch:livrable}

\begin{cle}
\textbf{Idée directrice.} Le chapitre~\ref{ch:identifiabilite} a montré que la direction de la
contagion n'est pas identifiable sur les sources agrégées. Cette limite n'est pas un point final :
elle se retourne en \textbf{spécification de reporting}, et cette spécification se \textbf{chiffre}
puis se \textbf{hiérarchise}. C'est la contribution réglementaire du travail, et elle ne dépend
d'aucun prior.
\end{cle}

\section{Ce qu'un registre devrait contenir}

Puisque les deux sources agrégées échouent pour des raisons \emph{opposées} (l'une manque de
taxonomie et de volume informatif, l'autre de résolution temporelle), la donnée qui identifierait
$W$ se déduit par complémentarité. Un registre d'incidents \textsc{tic} l'identifiant devrait :
\begin{enumerate}[itemsep=2pt]
\item horodater la \textbf{survenance} (et non la seule déclaration) au \textbf{mois ou à la
      semaine}, afin de porter une antériorité que le pas annuel efface ;
\item catégoriser par \textbf{domaine de contrôle} (les piliers DORA), et non par vecteur
      d'attaque, afin qu'une transition soit interprétable comme un lien $k \to j$ ;
\item consolider une \textbf{cause commune} comme \emph{un} événement à $N$ victimes, pour ne pas
      confondre un lot de déclaration avec une contagion séquentielle.
\end{enumerate}

Ces trois exigences ne sont pas hors d'atteinte : le registre d'information imposé par
l'\emph{implementing technical standard} du règlement d'exécution \citep{eu2024_2956} porte déjà
l'identification des prestataires, la chaîne de sous-traitance et la criticité des fonctions. Ce
qui manque relève du \emph{format} de l'horodatage et du \emph{choix de la taxonomie}, non d'une
collecte nouvelle.

\section{Ce que cette information vaut, en capital}

Une recommandation de reporting reste un v\oe u tant qu'elle n'est pas chiffrée. On la chiffre en
restreignant l'ensemble admissible du chapitre~\ref{chap:partielle} aux configurations compatibles
avec l'information acquise, puis en mesurant le rétrécissement de la bande de capital.

\begin{definition}[Valeur d'information d'une contrainte]\label{def:voi}
Soit $\mathcal{W}$ l'ensemble admissible et $\mathrm{SCR}:\mathcal{W}\to\R_+$ la fonctionnelle de
capital. La largeur de bande associée à un sous-ensemble $\mathcal{V}\subseteq\mathcal{W}$ est
\begin{equation}
  L(\mathcal{V}) \;=\; \sup_{W\in\mathcal{V}}\mathrm{SCR}(W)\;-\;\inf_{W\in\mathcal{V}}\mathrm{SCR}(W).
  \label{eq:largeur-bande}
\end{equation}
Une information $\mathcal{I}$ qui restreint l'ensemble à $\mathcal{W}_{\mathcal{I}}\subseteq
\mathcal{W}$ a pour \textbf{valeur} le rétrécissement relatif qu'elle produit :
\begin{equation}
  V(\mathcal{I}) \;=\; 1-\frac{L(\mathcal{W}_{\mathcal{I}})}{L(\mathcal{W})} \;\in\;[0,1].
  \label{eq:voi}
\end{equation}
\end{definition}

Cette définition a deux propriétés utiles. Elle est \textbf{monotone} : si
$\mathcal{I}'$ implique $\mathcal{I}$, alors $\mathcal{W}_{\mathcal{I}'}\subseteq
\mathcal{W}_{\mathcal{I}}$ et donc $V(\mathcal{I}')\ge V(\mathcal{I})$, une information plus
riche ne pouvant pas valoir moins. Et elle est \textbf{sans dimension}, donc indépendante du
niveau de capital, ce qui la rend comparable d'une source de sévérité à l'autre alors même que
ce niveau, lui, varie d'un facteur quarante.

Le corpus de sept rapports post-incident du chapitre~\ref{ch:identifiabilite} établit le
\emph{sens} de six des dix dépendances entre piliers. À lui seul, il fait tomber la largeur de la
bande de $\mathbf{66\,\%}$. C'est la valeur, en capital, de la lecture de sept documents publics,
et elle donne l'ordre de grandeur de ce qu'un registre correctement spécifié apporterait.

\begin{attention}
\textbf{Ce qu'aucune information de sens ne donnera.} Même les dix directions établies, une bande
subsisterait, imputable à la seule \textbf{amplitude} des coefficients : connaître le sens d'un
lien ne donne pas sa force. L'identification partielle n'est donc pas un état provisoire que la
donnée finirait par lever, elle est \textbf{structurelle}. C'est une raison de plus de lire le
capital en bande, et de ne pas promettre à un superviseur davantage que ce qu'un registre peut
tenir.
\end{attention}

\section{Quelle dépendance documenter en priorité}

L'apport le plus actionnable est ailleurs que dans le total. En fixant le sens d'\emph{une seule}
dépendance à la fois, on obtient un \textbf{classement} des données manquantes par valeur
d'information (figure~\ref{fig:voi}). L'écart est considérable : documenter la seule dépendance
entre gouvernance et gestion des incidents lèverait à elle seule $\mathbf{28{,}9\,\%}$ de
l'ambiguïté de capital, tandis que documenter le lien entre tests et tiers n'en lèverait
\emph{rien}.

Deux conséquences pratiques. D'abord, un superviseur qui ne pourrait exiger qu'une seule
dépendance sait désormais laquelle demander, et une entité qui ne pourrait instrumenter qu'un seul
lien sait lequel instrumenter. Ensuite, un fait qui oriente la recommandation : \textbf{la
dépendance la plus précieuse est précisément celle que le corpus documentaire ne couvre pas}. La
lecture de rapports publics a livré ce qu'elle pouvait livrer ; le reste exige la donnée d'entité.

\begin{figure}[htbp]
\centering
\includegraphics[width=\linewidth]{Z18_valeur_information.png}
\caption{La valeur de l'information, chiffrée et hiérarchisée. (a)~l'information rétrécit la
bande de capital : $66\,\%$ acquis par la lecture de sept rapports, et une bande d'amplitude qui
subsisterait même les dix sens connus. (b)~quelle dépendance documenter en priorité ; en bleu,
celles que le corpus couvre déjà. Source : script~\texttt{55}.}
\label{fig:voi}
\end{figure}

\begin{cle}
\textbf{La limite devenue contribution.} Un mémoire qui se contenterait de constater qu'il ne
peut pas calibrer $W$ laisserait le lecteur sans suite. En spécifiant la donnée qui manque, en
chiffrant ce qu'elle vaut et en la hiérarchisant, on transforme une frontière d'identification en
\textbf{recommandation opposable}, adressée au cadre de reporting que DORA met précisément en
place. C'est le seul endroit du mémoire où une limite produit un livrable.
\end{cle}

% =====================================================================
